\chapter{CONCLUSÃO}
\label{cap:conclusao}

Este trabalho investigou a integração de previsões quantitativas e informações textuais em um processo de alocação de ações brasileiras baseado no modelo de Black--Litterman. A avaliação permitiu examinar tanto o funcionamento da arquitetura proposta quanto o desempenho financeiro das carteiras resultantes. As conclusões apresentadas a seguir distinguem a implementação do fluxo de tratamento das opiniões da evidência empírica sobre sua capacidade de produzir resultados superiores aos referenciais.

\section{Síntese do Trabalho e Atendimento dos Objetivos}
\label{sec:conclusao_sintese}

O objetivo de desenvolver uma abordagem híbrida de otimização de portfólios foi atendido no âmbito da implementação e da avaliação experimental. Foi estruturado um fluxo que reúne séries históricas de mercado, atributos quantitativos, previsões do XGBoost, classificação de sentimento de manchetes e atualização de retornos esperados pelo Black--Litterman. Essas etapas foram conectadas à geração de pesos e à apuração de resultados financeiros, permitindo acompanhar o percurso entre as informações de entrada e as decisões de investimento simuladas.

No componente quantitativo, o trabalho considerou modelos individuais por ativo e previsões em diferentes horizontes. A escolha entre cinco, dez e vinte e um pregões permitiu representar escalas temporais distintas, em vez de impor um único horizonte a todo o universo. Na rodada consolidada, essa escolha foi realizada com base na acurácia direcional da validação de 2024 e condicionada à disponibilidade integral dos respectivos alvos antes de 2025. A utilização de previsões produzidas pelo procedimento \emph{walk-forward} possibilitou avaliar o comportamento do estimador em segmentos posteriores às suas janelas de treinamento, respeitadas as limitações da seleção original dos artefatos.

No componente qualitativo, as manchetes foram convertidas em sinais de sentimento e associadas às empresas do universo. A combinação com as previsões quantitativas originou opiniões de retorno, cuja influência foi controlada pela cobertura de notícias e pela incerteza atribuída a cada opinião. Assim, o sistema não se limitou a reunir duas fontes de informação: estabeleceu condições explícitas para sua participação na atualização das expectativas utilizadas pela carteira.

O fluxo de governança das opiniões foi implementado por meio da formação do vetor de views, da estimativa da matriz de incerteza e do mecanismo de \emph{hard prior}. Na configuração analisada, a presença de pelo menos três notícias no dia elegível habilita a opinião híbrida; abaixo desse limiar, sua confiança é reduzida. O uso de erros de previsões antigas somente após o encerramento de seus horizontes também torna explícita a distinção entre a informação disponível para decidir e aquela utilizada posteriormente para avaliar. O funcionamento dessas regras foi verificado na execução e na auditoria dos sinais, sem que isso equivalha a demonstrar sua calibração financeira ótima.

A avaliação do portfólio ampliou o alcance da comparação inicialmente centrada em uma carteira de tamanho fixo. Foram examinadas cardinalidades de um a vinte ativos em dois cenários: subconjuntos comuns aos modelos e seleção própria por estratégia. Essa separação permitiu investigar o desempenho da alocação sob oportunidades equivalentes e, complementarmente, o resultado do processo completo de seleção e distribuição de pesos. Ao todo, foram executadas 1.142 simulações individuais, com custos de transação e defasagem de execução, além de registros das composições e das trajetórias de capital.

Dessa forma, o trabalho entregou uma estrutura operacional e auditável para comparar o modelo híbrido com Markowitz e com uma referência de renda fixa. O atendimento desse objetivo não pressupõe que o sistema deva superar os referenciais em qualquer configuração. A contribuição experimental inclui identificar situações favoráveis, condições de desempenho inferior e limites da própria evidência, preservando a distinção entre viabilidade técnica e superioridade financeira.

\section{Principais Conclusões e Contribuições}
\label{sec:conclusao_contribuicoes}

Na avaliação exploratória de 2025 com o universo completo de vinte ativos, a carteira Black--Litterman apresentou retorno líquido acumulado de 28,21\%, contra 19,10\% de Markowitz e 14,87\% da referência de Selic constante de 15\% ao ano nas mesmas datas. Nesse recorte, o Sharpe calculado sobre o excesso de retorno foi de aproximadamente 0,73 para o modelo híbrido e 0,29 para Markowitz. O drawdown máximo foi de $-8{,}46\%$ e $-12{,}78\%$, respectivamente. A comparação indica uma combinação favorável de retorno e risco para a configuração examinada, sem estender essa conclusão automaticamente a outros períodos.

O estudo por cardinalidade mostrou que o desempenho depende tanto da composição quanto das restrições impostas à carteira. Nos subconjuntos comuns, os modelos coincidem em uma posição única e também em duas posições, pois o teto de 50\% obriga pesos-alvo iguais nesse segundo caso. Nas demais cardinalidades avaliadas, o retorno mediano do Black--Litterman foi superior ao de Markowitz. Essa evidência descreve os subconjuntos utilizados e não representa uma probabilidade de superioridade em um novo ano, uma vez que as carteiras compartilham ativos, datas e choques de mercado.

Na seleção própria, o Black--Litterman superou Markowitz em retorno em dezenove das vinte cardinalidades, mas apresentou perda de 10,37\% quando restrito a uma única ação. Esse contraste demonstra que resultados favoráveis em carteiras diversificadas podem coexistir com decisões de seleção malsucedidas. Também evidencia a utilidade de separar seleção e alocação: quando as ações elegíveis diferem, a diferença financeira não pode ser atribuída exclusivamente ao método de distribuição dos pesos.

Os resultados quantitativos recomendam cautela na interpretação dos ganhos da carteira. A acurácia direcional média entre ativos foi de 52,44\% em 2025, e dezenove dos vinte ativos apresentaram coeficiente de determinação negativo. Portanto, o desempenho financeiro observado não corresponde a uma demonstração de elevada precisão na previsão das magnitudes dos retornos. A carteira depende conjuntamente do prior, da covariância, das restrições, das views e de sua confiança, o que impede atribuir seu resultado a uma única fonte de informação.

A cobertura efetivamente utilizada reforça essa limitação de atribuição. Com gate de três notícias, apenas seis dos 240 pares ativo-data considerados nas decisões mensais habilitaram as views, distribuídos por três meses. O prior permaneceu relevante nas demais situações. A implementação confirma que o mecanismo regula a participação das opiniões, mas não demonstra, isoladamente, quanto o sentimento ou o XGBoost acrescentaram ao desempenho. Essa distinção constitui uma conclusão central: verificar o funcionamento de uma regra de governança não substitui medir a contribuição incremental das informações por ela governadas.

O alcance dos resultados também é limitado pela escolha retrospectiva do gate. O limiar de três notícias foi adotado após a comparação de valores entre um e cinco no próprio ano de 2025. Consequentemente, a vantagem reportada pertence a uma análise exploratória, e não a um teste temporal independente desse hiperparâmetro. Somam-se a isso a seleção original dos estimadores quantitativos, que ainda exige validação temporal integral, o universo fixo de ações e as simplificações de execução e custos. O prior equiponderado da rodada também deve ser distinguido da especificação ponderada por capitalização apresentada na arquitetura. Tais limitações não anulam o experimento, mas restringem as afirmações que podem ser sustentadas por ele.

A principal contribuição do trabalho é, assim, a integração de aprendizado de máquina, processamento de informações textuais e controle explícito de confiança em um fluxo de alocação verificável para o mercado brasileiro. A comparação por cardinalidade e por regra de seleção acrescenta uma forma de avaliar a dependência dos resultados em relação às oportunidades de investimento. O conjunto dessas contribuições oferece uma base para aperfeiçoamentos e testes posteriores, sem pressupor que a integração das fontes seja suficiente, por si só, para assegurar rentabilidade superior.

\section{Trabalhos Futuros}
\label{sec:conclusao_trabalhos_futuros}

A prioridade metodológica é realizar uma avaliação temporal independente de todas as decisões de seleção. Isso inclui a escolha entre classificador e regressor, o horizonte por ativo, o gate de notícias, os parâmetros de confiança e a cardinalidade operacional. Recomenda-se um protocolo de validação aninhada, no qual cada escolha seja determinada apenas por dados anteriores à janela avaliada, com separações compatíveis com a duração dos alvos prospectivos. Como 2025 já foi examinado na escolha do limiar, a confirmação da configuração requer um período que não tenha participado dessa decisão.

Outra etapa necessária é comparar versões controladas do sistema: Black--Litterman apenas com prior, com previsões quantitativas, com sentimento e com a combinação completa. Mantidas as mesmas restrições, datas e condições de execução, essa análise permitiria investigar se os ganhos decorrem da informação adicional, da regularização pelo prior ou de sua interação. Também seria pertinente avaliar a estabilidade dos resultados em diferentes janelas históricas, sem tratar os subconjuntos de uma mesma janela como observações independentes.

Na dimensão informacional, sugere-se ampliar o tratamento das notícias para além das manchetes, incorporando o corpo integral das reportagens, quando disponível, e seus horários de publicação. O objetivo seria investigar se informações contextuais melhoram a classificação de sentimento, sem presumir que textos mais extensos produzam sinais superiores. Esse avanço deve considerar duplicidade entre veículos, repercussão de um mesmo evento e associação a diferentes classes de ações da mesma companhia, evitando confundir repetição de conteúdo com aumento de evidência independente.

A relação entre notícia, confiança e frequência de decisão também merece aprofundamento. A comparação entre uma janela curta de notícias e uma agregação retrospectiva com decaimento temporal poderia reduzir a dependência de um único dia anterior ao rebalanceamento. Em paralelo, a matriz de incerteza pode ser calibrada com base no erro preditivo efetivamente observado e na qualidade da informação textual. Essas alternativas devem ser selecionadas em validação anterior e confrontadas com a especificação atual, incluindo análises homogêneas das frequências diária, semanal e mensal.

Para aproximar a avaliação das condições operacionais, recomenda-se a implementação de \emph{paper trading}. O registro das informações disponíveis, dos sinais, das ordens simuladas e dos preços de execução permitiria observar atrasos de coleta, disponibilidade dos dados, variação dos custos e diferenças entre o fechamento de referência e o preço negociável. A etapa também permitiria verificar se a rastreabilidade obtida no ambiente histórico se mantém em um fluxo de atualização contínua, antes de qualquer avaliação com recursos financeiros reais.

Como extensão posterior, a incorporação de posições vendidas ampliaria o conjunto de decisões possíveis, mas exigiria reformular as restrições e o modelo de custos. A avaliação deveria contemplar disponibilidade e custo de aluguel, limites de exposição, garantias e procedimentos de controle de risco. Também seriam relevantes dados históricos de composição do universo e de capitalização, além de uma modelagem mais detalhada de liquidez, tributação e impacto de mercado. Essas extensões devem ser examinadas pelo seu efeito conjunto sobre retorno, risco e viabilidade operacional.

Por fim, o desenvolvimento de novas versões deve preservar a separação entre construção, validação e confirmação dos resultados. A continuidade do trabalho não depende apenas de acrescentar variáveis ou tornar o modelo mais complexo, mas de produzir evidências que permitam identificar quais componentes contribuem de forma estável e em quais condições. Esse é o caminho para transformar a arquitetura experimental desenvolvida em uma ferramenta de apoio à decisão com alcance metodológico e operacional mais bem delimitado.
